\conclusion

В работе решена задача определения области применимости колоночного формата Vortex в
lakehouse-системе Apache Paimon, разработки модуля его интеграции и сравнительного
анализа форматов хранения. Поставленные задачи выполнены.

\begin{enumerate}
    \item Проведён аналитический обзор эволюции систем хранения аналитических данных
        --- от хранилищ данных и Hive Metastore к архитектуре lakehouse и открытым
        табличным форматам (Iceberg, Delta Lake, Hudi, Paimon); обоснован выбор
        Apache Paimon как целевой системы (LSM-дерево, ориентация на потоковую запись
        и таблицы с первичным ключом).
    \item Выполнен обзор колоночных файловых форматов (Parquet, ORC, Vortex, Lance);
        обосновано исключение строчного формата Avro из сравнения; разобраны
        алгоритмы кодирования, сжатия и пропуска данных --- словарное и RLE-кодирование,
        бит-упаковка, FOR и дельта-кодирование, FastLanes, ALP, FSST, статистики и
        фильтры Блума, проталкивание предикатов; выявлено архитектурное ограничение
        интерфейса форматов Apache Paimon --- отсутствие проталкивания агрегатов.
    \item Изучена архитектура Apache Paimon --- метаданные (снимки, манифесты),
        разбиение и бакеты, LSM-дерево таблицы с первичным ключом, слияние, путь
        чтения через итератор слияния, точка подключения файлового формата, интеграция
        с Apache Flink.
    \item Разработан и интегрирован в Apache Paimon модуль поддержки формата Vortex:
        реализация интерфейса \texttt{FileFormat} (фабрики записи и чтения),
        \texttt{VortexRecordsReader}, конвертер предикатов
        \texttt{VortexPredicateConverter}, слой доступа к нативной библиотеке через
        Java Native Interface с передачей данных без копирования через Arrow C Data
        Interface, формирование статистик записанных файлов, поддержка проекции
        столбцов и многопоточного чтения.
    \item Проведено функциональное тестирование артефакта локально в контейнерах:
        запись и чтение таблиц с первичным ключом и таблиц только для добавления,
        слияние версий, эволюция числа бакетов, проекция столбцов и проталкивание
        предикатов, граничные случаи кодеков; выявлен и устранён ряд проблем сборки и
        конфигурации.
    \item Развёрнуто на кластере из пяти узлов окружение нагрузочного тестирования
        средствами Ansible (Apache Flink, Apache Paimon с модулем
        \texttt{paimon-vortex}, объектное хранилище); автоматизированы развёртывание,
        обновление конфигурации, запуск заданий, сбор метрик, управление жизненным
        циклом кластера.
    \item Проведены две серии экспериментов --- batch-аналитика на наборе TPC-DS
        масштаба \SI{100}{\giga\byte} и near-real-time-сценарий с потоковой записью и
        параллельными аналитическими запросами; собраны и статистически обработаны
        метрики; результаты проанализированы по типам запросов.
    \item Сформулированы рекомендации по выбору файлового формата хранения для
        различных классов нагрузок в Apache Paimon.
\end{enumerate}

Гипотезы работы подтверждены. На batch-аналитике с селективными предикатами и
крупными файлами Vortex --- быстрейший из четырёх форматов (быстрее всех на
\num{90} запросах из \num{104}; по суммарному времени опережает Parquet на
\SI{25}{\percent}, ORC --- на \SI{32}{\percent}, пиковое ускорение --- до
\SI{83}{\percent}; Lance самый медленный) --- это основная область применимости
Vortex. На потоковой
near-real-time-нагрузке Vortex эквивалентен Parquet и ORC по пропускной способности
записи; по типичной задержке чтения (медиане) Vortex --- быстрейший формат в обеих
фазах: выигрывает девять запросов из тринадцати как в фазе FINAL, так и в фазе
DURING, особенно --- на запросах с селективным предикатом, диапазонным фильтром по
строке-префиксу и в многошаговой аналитике. На простой группировке и многомерной
агрегации без селективного предиката Vortex незначительно (в пределах единиц
процентов) уступает конкурентам --- следствие отсутствия проталкивания агрегатов в
интерфейсе форматов Apache Paimon. По хвостам задержки в фазе DURING ORC и Vortex
показывают узкие распределения (среднее почти равно медиане), Parquet --- более
широкое распределение с эпизодическими выбросами в десятки секунд в моменты сброса
данных по чекпойнту: это согласуется с архитектурным предположением о выигрыше
Vortex за счёт независимого нативного пула соединений с объектным хранилищем по
сравнению с общим пулом слоя совместимости Hadoop S3A, через который ходят
читатели Parquet и ORC. ORC достигает аналогично узких хвостов иначе ---
меньшим числом обращений к хранилищу за один скан за счёт крупных полос (stripes)
и гранулярного пропуска данных по индексным группам. Между фазами FINAL
и DURING медианы ORC и Parquet практически не меняются, медиана Vortex в фазе
FINAL заметно лучше его же медианы в DURING ($\approx 22\,\%$ улучшения) ---
эффект уплотнения данных слиянием, при котором специализированные кодеки Vortex
(прежде всего FSST на строках с регулярной структурой) работают в полной
плотности.

Практическая значимость работы: расширен перечень поддерживаемых Apache Paimon
форматов хранения; получены количественные результаты и рекомендации, применимые при
проектировании витрин данных и настройке lakehouse-хранилищ; систематизированы
проблемы развёртывания и эксплуатации Apache Paimon с нативным файловым форматом на
кластере с объектным хранилищем.

Направления дальнейшей работы: проталкивание агрегатов в нативный SIMD-слой Vortex
(устранение выявленного архитектурного ограничения и потенциальный вклад в
экосистему Apache Paimon); поддержка дополнительных операций проталкивания и типов
данных в конвертере предикатов; исследование поведения форматов при ещё более высоком
параллелизме и иных профилях нагрузки; доработка интеграции формата Lance для
таблиц с первичным ключом.
